%=======================================================================
%  Gercek Zamanli Fizyolojik Stres Tespitinde Aciklanabilirligin
%  Hesaplama Maliyeti -- ic teknik rapor (IEEE konferans bicimi)
%
%  Derleme:  pdflatex main  ->  pdflatex main   (iki kez)
%=======================================================================
\documentclass[conference]{IEEEtran}
\IEEEoverridecommandlockouts

\usepackage[T1]{fontenc}
\usepackage[utf8]{inputenc}
\usepackage[turkish]{babel}
\usepackage{amsmath,amssymb}
\usepackage{graphicx}
\usepackage{booktabs}
\usepackage{multirow}
\usepackage{array}
\usepackage{url}
\usepackage{tikz}
\usepackage{pgfplots}
\pgfplotsset{compat=1.17}
\usetikzlibrary{arrows.meta,positioning,fit}
\usepackage{xspace}
\usepackage[hidelinks]{hyperref}

\pgfkeys{/pgf/number format/.cd,use comma,1000 sep={}}

% Turkce babel '=', ':' ve '!' karakterlerini etkin hale getirir; matematik
% kipinde ve tablo kurulumunda bu sorun cikarir. Kisayollar kapatilir,
% tireleme ve "Sekil / TABLO / Kaynaklar" adlandirmasi korunur.
\newcommand{\kapatKisayol}{\shorthandoff{=:!}}

\newcommand{\ms}{\,ms\xspace}

\begin{document}
\kapatKisayol

\title{Gerçek Zamanlı Fizyolojik Stres Tespitinde\\
Açıklanabilirliğin Hesaplama Maliyeti:\\
Uç ve Sis Yerleşiminde Uçtan Uca Ölçüm}

\author{%
\IEEEauthorblockN{Sena Cindioğlu}
\IEEEauthorblockA{Bilgisayar Mühendisliği Bölümü\\
Yıldız Teknik Üniversitesi\\
İstanbul, Türkiye}
\and
\IEEEauthorblockN{Meryem Eldemery}
\IEEEauthorblockA{Bilgisayar Mühendisliği Bölümü\\
Yıldız Teknik Üniversitesi\\
İstanbul, Türkiye}
\and
\IEEEauthorblockN{Hafiza İrem Türkmen Çilingir}
\IEEEauthorblockA{Bilgisayar Mühendisliği Bölümü\\
Yıldız Teknik Üniversitesi\\
İstanbul, Türkiye}
}

\maketitle

\begin{abstract}
Fizyolojik sinyallerden stres tespiti literatüründe sınıflandırma başarımı
doyuma yaklaşmış, açıklanabilir yapay zekâ (XAI) yöntemleri de yaygınlaşmıştır.
Buna karşılık açıklama üretmenin \emph{hesaplama maliyeti} neredeyse hiç
ölçülmemektedir: yaygın XAI değerlendirme çerçeveleri çalışma süresini kapsam
dışı bırakmakta, WESAD üzerinde XAI uygulayan çalışmalar ise pencere başına
açıklama gecikmesi bildirmemektedir. Bu çalışma, bitirme projesinde üretilmiş
ve dondurulmuş bir stres sınıflandırıcısını gerçek zamanlı bir yeniden oynatma
(replay) sistemine yerleştirerek bu boşluğu doldurur. Sistem, sensör istemcisi
ile çıkarım sunucusunu gerçek bir TCP bağlantısıyla ayırır; gecikme, kuyruk
gecikmesi, çalışma süresi, yanıt süresi, seğirme, bant genişliği ve kaynak
kullanımı olmak üzere yedi metriği yüksek çözünürlüklü saatlerle, periyodik
saat yeniden senkronizasyonu altında ölçer. Bölge tabanlı LIME ve SHAP
GradientExplainer açıklayıcıları sunucunun içine, tahminle aynı iş parçacığına
yerleştirilmiş; hiperparametreleri sadakat, monotonluk, duyarlılık ve
karmaşıklık ölçütleriyle birlikte taranmıştır. Ayrı bir fiziksel makinede,
Wi-Fi üzerinden yapılan uçtan uca ölçümlerde açıklayıcı kapalıyken yanıt süresi
medyanı 10{,}6\ms iken, en pahalı yapılandırmada 1456{,}2\ms'ye, yani
\textbf{137 katına} çıkmaktadır; sunucu bellek ayak izi 210\,MB'den 464\,MB'ye
yükselmektedir. Ucuz çalışma noktaları (SHAP 442{,}6\ms, LIME 645{,}2\ms)
literatürdeki bir saniyelik gerçek zamanlılık sınırının altında kalmaktadır. Bu
sayılar, güven eşikli seçici açıklama mekanizmasının niçin gerekli olduğunu
niceliksel olarak ortaya koymaktadır.
\end{abstract}

\begin{IEEEkeywords}
açıklanabilir yapay zekâ, stres tespiti, WESAD, sis bilişim, gerçek zamanlı
sistemler, hesaplama maliyeti, SHAP, LIME
\end{IEEEkeywords}

%=======================================================================
\section{Giriş}
%=======================================================================

Giyilebilir sensörlerden sürekli stres izleme, sağlık ve insan--bilgisayar
etkileşimi uygulamaları için olgunlaşmış bir problemdir. WESAD \cite{schmidt2018}
gibi ölçüt veri kümeleri üzerinde \%90'ın üzerinde doğruluk rutin biçimde
bildirilmekte \cite{kasnesis2025}, model kararlarının şeffaflığı için de SHAP
\cite{lundberg2017}, LIME \cite{ribeiro2016} ve Integrated Gradients
\cite{sundararajan2017} gibi post-hoc açıklayıcılar yaygın olarak
kullanılmaktadır.

Ancak bu iki gelişme birbirinden bağımsız ilerlemiştir. Sınıflandırma başarımını
raporlayan çalışmalar açıklamanın maliyetini, açıklama üreten çalışmalar da
sistemin zamanlamasını ölçmemektedir. Nauta ve ark.'nın Co-12 değerlendirme
çerçevesi bu dışlamayı açıkça yapar: \emph{``çalışma süresi gibi, açıklama
kalitesini doğrudan etkilemeyen ölçütler hariç tutulmuştur''} \cite{nauta2023}.
Quantus, OpenXAI ve XAI-Bench gibi kıyaslama araçları da süre sütunu içermez.
Sonuç olarak ``açıklama maliyeti'' için ne standart bir metrik, ne standart bir
birim, ne de liderlik tablolarında bir yer vardır.

Bu boşluk çevrimdışı bir çalışmada zararsızdır; gerçek zamanlı bir sistemde ise
belirleyicidir. Tek bir açıklama, pencere güncelleme aralığından uzun sürüyorsa
sistem kuyruk biriktirir ve gerçek zamanlı olmaktan çıkar. Franco de la Peña ve
ark. bu sınırı açıkça koyar: \emph{``pencere başına bir saniyeden fazla süren bir
XAI yöntemi gerçek zamanlı dağıtıma uygun değildir''} \cite{shats2025}.

Bu rapor, bir bitirme projesi kapsamında üretilmiş ve dondurulmuş bir stres
sınıflandırıcısını temel alarak bu eksiği kapatan ölçüm çalışmasını
belgelemektedir. Katkılar şunlardır:

\begin{enumerate}
\item \textbf{Sızıntısız çevrimiçi değerlendirme.} Çevrimdışı eğitimde kullanılan
denek bazlı z-skor normalizasyonunun nedensel (yalnızca geçmişi gören) karşılığı
uygulanmış ve doğruluk üzerindeki etkisi ölçülmüştür.
\item \textbf{Aşama bazlı gerçek zamanlı profil.} Yeniden örnekleme, ön işleme,
tensör hazırlığı, çıkarım ve son işleme aşamaları ayrı ayrı, tekrarlı ve
dönüşümlü koşularla, koşu içi ve koşu arası dağılımlar ayrıştırılarak
ölçülmüştür.
\item \textbf{Yedi metrikli taşıma katmanı ölçümü.} İstemci ve sunucu gerçek bir
TCP bağlantısıyla ayrılmış, iki bağımsız yüksek çözünürlüklü saat arasında
periyodik yeniden senkronizasyon ve doğrusal ara değerleme uygulanmıştır.
\item \textbf{Açıklama maliyetinin taşıma metrikleriyle aynı bütçede
raporlanması.} Bilgimiz dahilinde, pencere başına açıklama gecikmesini kuyruk
gecikmesi, seğirme, bant genişliği ve sunucu kaynak kullanımıyla birlikte veren
bir çalışma yoktur.
\item \textbf{Açıklayıcı hiperparametre taraması.} LIME örnek sayısı, bölge
sayısı ve çekirdek genişliği ile SHAP arka plan boyutu, örnek sayısı ve arka plan
seçim yöntemi; süre \emph{ve} dört kalite ölçütü birlikte verilerek taranmıştır.
\end{enumerate}

Güven eşikli seçici açıklama mekanizması sisteme entegre edilmiş ve iki uç durum
(hiç açıklama / her pencerede açıklama) uçtan uca ölçülmüştür. Ara eşiklerin
taraması ve güven kalibrasyonu devam eden iştir; Bölüm~\ref{sec:devam}'de açıkça
belirtilmiştir.

%=======================================================================
\section{Önceki Aşama: Çevrimdışı Sınıflandırma ve Açıklanabilirlik}
\label{sec:onceki}
%=======================================================================

Bu bölüm, üzerine inşa edilen bitirme projesinin \cite{bitirme2026} kısa
özetidir.

WESAD veri kümesinde \cite{schmidt2018} göğüs (RespiBAN) ve bilek (Empatica E4)
sensörlerinden gelen EDA, EKG, EMG, solunum, deri sıcaklığı, kan hacmi nabzı ve
üç eksenli ivmeölçer sinyalleri zaman pencerelerine bölünerek stres / stres dışı
olarak sınıflandırılmıştır. Üç mimari karşılaştırılmıştır: bir evrişimsel sinir
ağı (CNN), hibrit bir CNN--LSTM ve Transformer tabanlı TranSensFusers
\cite{kasnesis2025}. Sınıf dengesizliği odak kaybı (focal loss)
\cite{lin2017} ile ele alınmış, değerlendirme denek bağımsız
Leave-One-Subject-Out protokolüyle yapılmıştır. Pencere boyutu her mimari için
ayrıca taranmış; en iyi sonuç 10 saniyelik pencerelerle eğitilen TranSensFusers
modelinden elde edilerek \%91{,}13 doğruluk ve \%90{,}40 F1-skoru elde edilmiştir
(Tablo~\ref{tab:onceki-model}).

\begin{table}[t]
\caption{Bitirme projesindeki en iyi yapılandırmalar (özet)}
\label{tab:onceki-model}
\centering
\begin{tabular}{lccc}
\toprule
\textbf{Model} & \textbf{Pencere} & \textbf{Doğruluk (\%)} & \textbf{F1 (\%)}\\
\midrule
CNN            & 30\,s & 87{,}63 & 86{,}29\\
CNN--LSTM      & 30\,s & 86{,}08 & 83{,}66\\
TranSensFusers & 10\,s & \textbf{91{,}13} & \textbf{90{,}40}\\
\bottomrule
\end{tabular}
\end{table}

En iyi modele SHAP, LIME ve Integrated Gradients (IG) uygulanmış; açıklamalar
sadakat, monotonluk, duyarlılık ve karmaşıklık olmak üzere dört nicel ölçütle
değerlendirilmiştir (Tablo~\ref{tab:onceki-xai}). SHAP en dengeli ve en kararlı
açıklamaları üretmiş, IG tamamlayıcı gradyan tabanlı bilgi sağlamış, LIME ise
seyrek fakat daha az kararlı açıklamalar vermiştir. Kanal düzeyinde atıf
analizinde elektrodermal aktivite ve kardiyak sinyaller baskın katkı sağlayarak
fizyolojik beklentiyi doğrulamıştır.

\begin{table}[t]
\caption{Bitirme projesindeki XAI değerlendirmesi (TranSensFusers, 10\,s pencere)}
\label{tab:onceki-xai}
\centering
\small
\setlength{\tabcolsep}{4pt}
\begin{tabular}{lcccc}
\toprule
\textbf{Ölçüt} & \textbf{SHAP} & \textbf{LIME} & \textbf{IG} & \textbf{En iyi}\\
\midrule
Sadakat (ekleme)   & 0{,}4048 & 0{,}4282 & \textbf{0{,}4402} & \multirow{2}{*}{IG}\\
Sadakat (silme)    & 0{,}2613 & 0{,}2153 & \textbf{0{,}2055} & \\
Monotonluk ihlali  & \textbf{0{,}4047} & 0{,}5079 & 0{,}5300 & SHAP\\
Duyarlılık         & \textbf{0{,}0038} & 0{,}0116 & 0{,}0053 & SHAP\\
Karmaşıklık        & 0{,}4674 & \textbf{0{,}0066} & 0{,}4585 & LIME\\
\bottomrule
\end{tabular}
\\[2pt]
\footnotesize Ekleme eğrisinde yüksek, diğer üç ölçütte düşük değer iyidir.
\end{table}

\textbf{Bu çalışmanın çıkış noktası.} Bitirme projesinin bütün ölçümleri
çevrimdışıdır: tüm kaydın istatistiğiyle normalize edilmiş pencereler, toplu
çıkarım, ağ yok, süre ölçümü yok. Dolayısıyla ``model yanıtı 100\ms altı,
açıklama 2\,s altı'' biçimindeki gereksinimler \emph{iddia} edilmiş ama test
edilmemiştir. Bu rapor tam olarak bu gereksinimleri ölçülebilir hale getirir.

\textbf{Kapsam daraltması.} Gerçek zamanlı çalışmada tek bir model (TemporalCNN)
ve iki açıklayıcı (SHAP, LIME) kullanılmaktadır. Gerekçe, karşılaştırmanın sabit
bir hesap bütçesi üzerinde yapılabilmesidir: TranSensFusers ve CNN--LSTM'in
eklenmesi açıklama maliyetindeki farkı model maliyetindeki farkla karıştırırdı.
Integrated Gradients, SHAP GradientExplainer ile aynı gradyan tabanlı maliyet
sınıfında olduğu için kapsam dışı bırakılmıştır.

%=======================================================================
\section{İlgili Çalışmalar}
%=======================================================================

\subsection{Sis mimarisinde fizyolojik izleme}

Erzurumluoğlu ve ark.\ \cite{erzurumluoglu2025}, DREAMER veri kümesi üzerinde
duygu tanımayı bir sis (fog) mimarisine yerleştirmiş ve gecikme, kuyruk
gecikmesi, çalışma süresi, yanıt süresi, seğirme, bant genişliği ve kaynak
kullanımı olmak üzere yedi metriği raporlamıştır. Bu çalışma, sistem ölçüm
çerçevemizin doğrudan referansıdır; aynı yedi metrik, aynı tanımlarla
kullanılmıştır. Farklar ikidir: görev (stres, WESAD) ve eklenen boyut
(açıklama maliyeti).

\subsection{Açıklama maliyeti ölçümleri}

Açıklayıcıların çalışma süresini raporlayan çalışmalar azdır ve büyük ölçüde
fizyolojik alan dışındadır. Tablo~\ref{tab:lit-maliyet} doğrulanabilmiş
değerleri toplar. Değerlerin donanım ve girdi boyutuna aşırı duyarlı olduğuna
dikkat edilmelidir: aynı yöntem (KernelSHAP) bir çalışmada saniyeler, başka bir
çalışmada dakikalar sürmektedir.

\begin{table}[t]
\caption{Literatürde bildirilen açıklama süreleri}
\label{tab:lit-maliyet}
\centering
\small
\begin{tabular}{p{2.1cm}p{2.4cm}p{2.7cm}}
\toprule
\textbf{Çalışma} & \textbf{Ölçülen} & \textbf{Donanım / girdi}\\
\midrule
ShaTS \cite{shats2025} & Sınır: $>$1\,s/pencere gerçek zamanlı değil; KernelSHAP sık sık aşıyor & i7-14700K + RTX 4080; girdi $10\times69$\\
XAI-on-RAN \cite{basaran2025} & Taban 5{,}3\ms; SHAP ($m{=}16$) $\approx +15$\ms & A100; 5 boyutlu girdi\\
Taguta ve ark. \cite{taguta2026} & Koşullu SHAP $+195$\ms, $+\%28$ CPU, toplam $\approx$555\ms & Raspberry Pi 4, 24 günlük canlı dağıtım\\
Das ve ark. \cite{das2021} & LIME 5{,}91\,s; SHAP 695{,}2\,s & i7-6700K; 33 sensör $\times$ 30 adım LSTM\\
Warnecke ve ark. \cite{warnecke2020} & KernelSHAP 9{,}1\,s; LIME 0{,}03--0{,}74\,s & Xeon E5; MLP/CNN\\
FastSHAP \cite{jethani2022} & 1000 görüntü: FastSHAP 0{,}04\,dk; KernelSHAP 453--1089\,dk & V100 + 8 çekirdek\\
\bottomrule
\end{tabular}
\end{table}

Fizyolojik sinyal tarafında durum daha çarpıcıdır. Almadhor ve ark.
\cite{almadhor2025} tam olarak bizim veri kümesi ailemizi (WESAD ve DREAMER) ve
model ailemizi (1B-CNN) kullanmakta, üstelik Grad-CAM ve Integrated Gradients
uygulamaktadır; buna karşın çıkarım süresi, XAI süresi, bellek veya donanım
bilgisi vermemektedir. Benzer biçimde El Rashidy ve ark.\ \cite{elrashidy2025}
sis mimarisinde LIME ve Grad-CAM ile duygu tanıma yapmakta, yine de çıkarım
başına gecikme ya da açıklama maliyeti raporlamamaktadır.

\subsection{Seçici ve koşullu açıklama}

Hesabı güvene göre dağıtma fikri yeni değildir. BranchyNet \cite{teerapittayanon2017}
ve Berestizshevsky ve Even \cite{berestizshevsky2018} softmax güvenine göre
çıkarımı erken keser. Açıklama tarafında Paes ve ark.\ \cite{paes2024} bütçeli
seçici açıklama önerir; kapı kararını ucuz açıklayıcının belirsizliği verir.
Artelt ve ark.\ \cite{artelt2023} yalnızca reddedilen örnekler için açıklama
üretir; gerekçeleri epistemiktir ve süre ölçmezler. Taguta ve ark.
\cite{taguta2026} SHAP'i koşullu tetikleyen, bilgimiz dahilinde donanım üzerinde
ölçülmüş tek sistem çalışmasıdır.

Dolayısıyla bu çalışma \emph{mekanizmayı} yeni diye sunmamaktadır. Katkı,
mekanizmanın sürekli gerçek zamanlı fizyolojik izleme bağlamında deneysel
olarak karakterize edilmesi ve açıklama maliyetinin taşıma katmanı metrikleriyle
aynı bütçede raporlanmasıdır.

\subsection{Boşluk}

Alan üç ayrık gruba bölünmüştür: (i) sistem metriklerini ölçen, XAI içermeyen
çalışmalar \cite{erzurumluoglu2025}; (ii) XAI maliyetini ölçen, ağ metriği
vermeyen çalışmalar \cite{shats2025,basaran2025}; (iii) fizyolojik sinyalde XAI
uygulayan, hiç maliyet vermeyen çalışmalar \cite{almadhor2025,elrashidy2025}. Bu
rapor üçünün kesişiminde konumlanmaktadır.\footnote{Bu bölümün hazırlığında
kullanılan iç literatür notlarında, doğrulanamayan veya kendi içinde tutarsız
sayılar içeren kaynaklar işaretlenmiş ve kaynakçaya alınmamıştır.}

%=======================================================================
\section{Sistem Tasarımı}
%=======================================================================

Sistemin genel akışı Şekil~\ref{fig:sistem}'de verilmiştir.

\begin{figure}[t]
\centering
\begin{tikzpicture}[
  node distance=3.2mm,
  blok/.style={draw,rounded corners=1.5pt,align=center,font=\scriptsize,
               minimum height=6.2mm,inner xsep=2pt},
  ok/.style={-{Latex[length=1.6mm]},thick}
]
\node[blok,fill=black!5,minimum width=17mm] (kayit) {WESAD kaydı\\(ham hız)};
\node[blok,fill=black!5,minimum width=17mm,below=of kayit] (tampon) {yuvarlanan tampon\\(500\ms tik)};
\node[blok,fill=black!5,minimum width=17mm,below=of tampon] (gonder) {paketle ve gönder\\(TCP)};
\node[draw,dashed,inner sep=2.0mm,fit=(kayit)(tampon)(gonder),
      label={[font=\scriptsize\bfseries]above:İSTEMCİ}] (istemci) {};

\node[blok,fill=black!10,minimum width=20mm,right=16mm of kayit] (kuyruk) {açık istek kuyruğu\\($t_{\mathrm{recv}}$ damgası)};
\node[blok,fill=black!10,minimum width=20mm,below=of kuyruk] (on) {yeniden örnekleme\\+ nedensel z-skor};
\node[blok,fill=black!10,minimum width=20mm,below=of on] (cikarim) {TemporalCNN\\çıkarım + softmax};
\node[blok,fill=black!20,minimum width=20mm,below=of cikarim] (kapi) {güven $<$ eşik ?};
\node[blok,fill=black!10,minimum width=20mm,below=of kapi] (xai) {LIME / SHAP\\açıklama};
\node[draw,dashed,inner sep=2.0mm,fit=(kuyruk)(on)(cikarim)(kapi)(xai),
      label={[font=\scriptsize\bfseries]above:SUNUCU (tek çalışan)}] (sunucu) {};

\draw[ok] (kayit) -- (tampon);
\draw[ok] (tampon) -- (gonder);
\draw[ok] (gonder.east) -- ++(4mm,0) |- (kuyruk.west);
\draw[ok] (kuyruk) -- (on);
\draw[ok] (on) -- (cikarim);
\draw[ok] (cikarim) -- (kapi);
\draw[ok] (kapi) -- node[right,font=\scriptsize]{evet} (xai);
\draw[ok] (kapi.east) -- ++(4mm,0) node[right,font=\scriptsize,xshift=-1mm]{hayır};
\end{tikzpicture}
\caption{Sistem akışı. Açıklama, tahminle aynı çalışanda ve çalışma süresinin
\emph{içinde} üretilir; ayrı bir iş parçacığına atılsaydı maliyeti ölçülemezdi.}
\label{fig:sistem}
\end{figure}

\subsection{Model ve pencereleme}

Sınıflandırıcı, çevrimdışı aşamada seçilip dondurulmuş bir TemporalCNN'dir:
iki evrişim bloğu (\texttt{[32, 128]} filtre, çekirdek boyu 9, her blokta
BatchNorm--ReLU--MaxPool--Dropout), ardından global ortalama havuzlama ve iki
tam bağlı katman. Toplam 49\,890 parametre içerir. Girdi 14 kanallı, 64\,Hz'e
indirgenmiş 10 saniyelik bir penceredir ($640\times14$); pencereler \%50
örtüşmeyle, yani 5 saniyede bir üretilir. Model 13 denekle eğitilmiş, S11 ve S14
tutulmuş (held-out) denek olarak ayrılmıştır; çevrimdışı ağırlıklı F1 değeri
0{,}9386'dır.

Kanal sırası eğitim ardışık düzeniyle birebir aynı olmak zorundadır (önce göğüs
sekiz kanal, sonra bilek altı kanal). Sıra bozulduğunda doğruluğun 1{,}00'den
0{,}52'ye düştüğü gözlenmiş, bu nedenle sıra kod içinde açıkça belgelenmiştir.

\subsection{Sızıntısız çevrimiçi normalizasyon}

Çevrimdışı eğitimde veri denek bazlı \texttt{StandardScaler} ile normalize
edilmişti. Gerçek zamanlı bir sistemde bu doğrudan kullanılamaz: bir deneğin tüm
kaydının istatistiğini bilmek geleceğe sızmaktır. İki sızıntısız alternatif
uygulanmış ve ölçülmüştür:

\begin{itemize}
\item \textbf{\texttt{running}:} Welford algoritmasıyla, yalnızca o ana kadar
görülmüş veriden nedensel ortalama ve standart sapma. İstatistik güncelleme
maliyeti her tikte ödenir ve dürüstlük gereği pencerenin ön işleme süresine
eklenir.
\item \textbf{\texttt{fixed}:} Yalnızca eğitim deneklerinden hesaplanmış tek bir
sabit ortalama ve standart sapma. Yeni kullanıcı hakkında hiçbir şey
bilinmeyen katı dağıtım senaryosuna karşılık gelir.
\end{itemize}

İki uzayın ürettiği pencereler ölçülmüştür: bağıl L2 farkı 0{,}77, standart
sapma 1{,}00'e karşı 1{,}25'tir. Yani nedensel ve denek bazlı normalizasyon
\emph{aynı girdi uzayını vermez}. Bu, XAI ölçümlerinin hangi uzayda yapılması
gerektiğini belirleyen kritik bulgudur (Bölüm~\ref{sec:xai-tarama}).

\subsection{İki ardışık düzen yolu}

Yeniden örneklemenin nerede yapıldığı bir tasarım kararıdır ve ölçüme doğrudan
girer:

\begin{itemize}
\item \textbf{\texttt{prepared}:} Tüm kayıt, yeniden oynatma başlamadan önce
64\,Hz'e indirgenir; yeniden örnekleme ölçümün \emph{dışında} kalır.
\item \textbf{\texttt{native}:} Her sensör kendi hızında akar (göğüs 700\,Hz,
bilek 32/64/4\,Hz), indirgeme sunucuda her pencerede yapılır ve \emph{ölçülür}.
\end{itemize}

Gerçek bir sistemde sensör ham hızında veri gönderir; dolayısıyla
\texttt{native} gerçekçi olan yoldur ve tüm ana sonuçlar bu yolda üretilmiştir.

\subsection{İstemci--sunucu ayrımı}

Erken aşamalarda istemci ve sunucu aynı süreçte çalışıyordu; bu, ağ kaynaklı
metriklerin tanım gereği sıfır olması demekti. \texttt{serving/} paketi ikisini
gerçek bir TCP bağlantısıyla ayırır: \texttt{protocol.py} tel formatını,
çerçevelemeyi ve saat senkronizasyonunu; \texttt{server.py} açık istek kuyruğunu
ve tek çalışanı; \texttt{client.py} ham hızda gönderen sensör istemcisini
tanımlar. Sunucu WESAD verisine ihtiyaç duymaz; veri tamamen istemcide kalır ve
sunucuya yalnızca ham sinyal parçaları akar.

\textbf{Açık istek kuyruğu neden gerekli.} İlk sürümde paket okunur okunmaz aynı
yerde işleniyordu; $t_{\mathrm{start}}$ ile $t_{\mathrm{recv}}$ bitişik olduğu
için kuyruk gecikmesi her zaman sıfıra yakın okunuyor ve bekleme süresi
gecikmenin içinde gizleniyordu. Şimdi okuyucu görev paketi alıp varış zamanını
damgalar ve ortak kuyruğa koyar; tek çalışan kuyruğu sırayla boşaltır. İki
istemciyle p99 kuyruk gecikmesi 0{,}37\ms'den 5{,}50\ms'ye çıkmıştır, yani
metrik artık gerçekten bekleme ölçmektedir.

\subsection{Saat katmanı ve kayma düzeltmesi}
\label{sec:saat}

İlk sürüm \texttt{time.time\_ns()} kullanıyordu. Windows'ta bu sayacın
çözünürlüğü ölçülmüş ve 15{,}6\ms çıkmıştır. Geri döngü (loopback) gecikmesi
0{,}1\ms mertebesinde olduğundan gecikme medyanı tam sıfır okuyordu; metrik
tamamen anlamsızdı. \texttt{perf\_counter\_ns} 100\,ns çözünürlük verir, ancak
referans noktası sürece özeldir. Bu nedenle bağlantı başında NTP mantığıyla saat
farkı ölçülür: on iki PING/PONG turu atılır ve en düşük gidiş-dönüş süresine
sahip tur seçilir.

Tek bir ölçüm uzun koşularda yetmez: iki kristal aynı hızda tiklemediği için
fark zamanla kayar. 87 dakikalık bir Wi-Fi koşusunda kayma, gecikme medyanını
\emph{negatife} düşürerek metriği geçersiz kılmıştır. Bu nedenle saat farkı
varsayılan olarak 60 saniyede bir yeniden ölçülür ve her paketin gecikmesi iki
ölçüm arasında doğrusal ara değerlemeyle düzeltilir. Ölçüm, son cevap alındıktan
sonraki sessiz anda dört turda yapılır. Özete toplam kayma ve kayma hızı (ppm)
yazılır; ayrıca gecikme medyanı
$(\text{yanıt} - \text{çalışma} - \text{kuyruk})/2$ ile karşılaştırılarak
tutarlılık denetimi basılır.

Gecikme belirsizliği gidiş-dönüş süresinin yarısı kadardır:
geri döngüde $\pm0{,}13$\ms, Ethernet'te $\pm0{,}25$--$1$\ms, Wi-Fi'da
$\pm1$--$10$\ms. Wi-Fi'da tek yönlü gecikmenin hata payı, ölçülen değerin
kendisiyle aynı mertebeye çıkabilir. Bundan \emph{etkilenmeyen} metrikler
çalışma süresi ve kuyruk gecikmesi (ikisi de sunucu saati) ile yanıt süresi ve
seğirmedir (ikisi de istemci saati).

\subsection{Ölçülen yedi metrik}

Tanımlar referans çalışmayla \cite{erzurumluoglu2025} aynıdır:

\begin{itemize}
\item \textbf{Gecikme (latency):} paketin istemciden sunucuya ulaşması; iki saat,
fark düzeltmesi gerektirir.
\item \textbf{Kuyruk gecikmesi:} çalışan meşgulken kuyrukta bekleme; sunucu saati.
\item \textbf{Çalışma süresi (execution):} ardışık düzenin çalışma süresi; sunucu
saati. Açıklama üretildiğinde bu sürenin içindedir.
\item \textbf{Yanıt süresi (response):} istemcinin gördüğü toplam süre; istemci
saati.
\item \textbf{Seğirme (jitter):} yanıt süresinin ardışık mutlak değişiminin
ortalaması.
\item \textbf{Bant genişliği:} gönderilen ve alınan bayt / süre.
\item \textbf{Kaynak kullanımı:} sunucu tarafı CPU yüzdesi ve yerleşik bellek
(RSS).
\end{itemize}

\subsection{Açıklayıcılar}

Açıklayıcılar sunucunun içinde, tek pencere üzerinde çalışacak biçimde
yazılmıştır. Ortak arayüz \texttt{explain(pencere)} çağrısıyla $(C,T)$ boyutlu
bir atıf haritası döndürür. Pencere yönü iki farklı gösterimde dolaştığı için
($(640,14)$ tampon yönü, $(14,640)$ model yönü) dönüşüm kanal sayısına bakarak
otomatik yapılır.

\textbf{Bölge tabanlı LIME.} Hazır \texttt{lime} paketi kullanılmamıştır: girdi
$14\times640 = 8960$ değerlik bir zaman serisidir ve her örneği ayrı değişken
saymak hem çok pahalı hem de yorumlanamazdır. Görüntüdeki ``süper piksel''in
buradaki karşılığı \emph{zaman dilimi $\times$ kanal} bölgeleridir; örneğin 20
dilim $\times$ 14 kanal $=$ 280 bölge. Bir bölge kapatıldığında o kanalın o
dilimi \emph{kanalın kendi ortalamasıyla} değiştirilir. Sıfırlamak fizyolojik
olarak anlamsız olurdu: normalize edilmiş sinyalde sıfır, o kanalın ortalama
seviyesi değil kaydın ortalamasıdır. Maskeler sabit tohumla örneklenir, ağırlık
$w = \sqrt{\exp(-d^2/\sigma^2)}$ üssel çekirdeğiyle verilir ($d$: kapalı bölge
oranı, $\sigma$: çekirdek genişliği) ve katsayılar Ridge regresyonuyla
çözülür. Maliyet, örnek sayısı $\times$ ileri geçiş olarak doğrusal büyür.

\textbf{SHAP GradientExplainer.} Arka plan kümesi bir kez üretilip diske
yazılır (100 pencere, k-ortalamalar, S2--S5 deneklerinden, sunucunun gördüğü
nedensel uzayda). Gerekçe ikilidir: sunucuda her açılışta yeniden kümeleme
saniyeler alır ve taramada kullanılan kümeyle birebir aynı sonucu vermez.
Tohum sabitlenmiştir; aksi halde aynı pencere iki kez açıklandığında farklı atıf
çıkar ve ``girdiye duyarlılık'' ölçütü yöntemin kendi rastgeleliğini de sayar.
Bu ölçülmüştür: tohum sabitlenince SHAP'in duyarlılık değeri 0{,}139'dan
0{,}075'e düşmüş, yani farkın neredeyse yarısının yöntemin kendi gürültüsü
olduğu görülmüştür.

\textbf{Isınma.} SHAP'in \emph{ilk} çağrısı ölçüm makinesinde 12{,}4 saniye
sürmüştür; sonraki çağrılar 0{,}85 saniyedir. Fark, grafik kurulumu ve önbellek
ısınmasıdır. Isıtılmazsa bu 12 saniye ilk açıklamada tek çalışanı bloklar ve tüm
kuyruk arkasında birikir. Bu nedenle açıklayıcı sunucu açılışında bir kez sahte
bir pencereyle çalıştırılır ve ilk çağrı süresi rapora yazılır.

\subsection{Açıklama kalitesi ölçütleri}

Dört ölçüt, bitirme projesindekiyle aynı tanımlarla hesaplanır:

\begin{itemize}
\item \textbf{Silme (deletion):} bölgeler en önemliden başlayarak kanal
ortalamasıyla kapatılır; eğri altındaki alanın \emph{düşük} olması sadakat
göstergesidir.
\item \textbf{Ekleme (insertion):} taban çizgisinden başlanarak en önemli
bölgeler eklenir; alanın \emph{yüksek} olması iyidir.
\item \textbf{Sadakat:} $\mathrm{AUC}_{\mathrm{ekleme}} -
\mathrm{AUC}_{\mathrm{silme}}$.
\item \textbf{Monotonluk:} ekleme eğrisinde azalmayan adımların oranı.
\item \textbf{Duyarlılık:} girdiye $0{,}05\sigma$ ölçeğinde Gauss gürültüsü
eklendiğinde atıfın normalize L2 değişimi; \emph{düşük} iyidir.
\item \textbf{Kendi-kararsızlığı:} aynı pencere yeniden açıklandığında atıfın
değişimi. Deterministik bir yöntemde sıfırdır. Bu ayrı raporlanmazsa ``kararsız
yöntem'' ile ``girdiye duyarlı yöntem'' birbirine karışır.
\end{itemize}

\subsection{Güven eşikli seçici açıklama}

Sunucu, softmax çıktısının maksimumunu güven olarak kullanır ve eşiğin altında
kalan pencerelerde açıklama üretir:

\begin{equation}
\text{açıkla} =
\begin{cases}
\text{hayır}, & \tau \le 0\\
\text{evet}, & \tau \ge 1\\
c < \tau, & \text{aksi halde}
\end{cases}
\end{equation}

\noindent burada $c$ güven, $\tau$ eşiktir. $\tau = 0$ ve $\tau = 1$ uç durumlar
kasıtlı olarak ayrı ele alınır; bunlar kontrol koşularıdır (XAI'siz taban ve en
pahalı durum). $\tau \ge 1$ karşılaştırması ayrıca ele alınmıştır çünkü softmax
çıktısı tam 1{,}0 okuyabildiğinde ``$c < 1{,}0$'' kuralı nadiren de olsa
pencere atlıyordu.

Açıklama, tahminle aynı çalışanda ve çalışma süresinin içinde üretilir. Ayrı bir
iş parçacığına atılsaydı maliyeti görünmez olurdu; deneyin bütün amacı o maliyeti
ölçmektir.

%=======================================================================
\section{Deney Düzeneği}
%=======================================================================

\subsection{Donanım ve ağ}

Üç yerleşim ölçülmüştür (Tablo~\ref{tab:donanim}). ``Uç'' terimi burada
\emph{ortak yerleşim} anlamındadır: istemci ve sunucu aynı makinede, ayrı
süreçlerde, geri döngü üzerinden haberleşir. Giyilebilir cihaz anlamında
kullanılmamaktadır.

\begin{table}[t]
\caption{Ölçüm ortamları}
\label{tab:donanim}
\centering
\small
\begin{tabular}{p{1.3cm}p{3.1cm}p{2.7cm}}
\toprule
\textbf{Yerleşim} & \textbf{Sunucu donanımı} & \textbf{Bağlantı}\\
\midrule
Geri döngü & İstemciyle aynı makine (Intel, 20 çekirdek, 16\,GB, Win 10) & 127.0.0.1\\
Uç (ortak) & Aynı makine, ayrı süreç & 127.0.0.1\\
Sis & Ayrı fiziksel makine (Intel, 12 çekirdek, 16\,GB, Win 11, PyTorch 2.14 CPU) & Wi-Fi, ev yönlendiricisi\\
\bottomrule
\end{tabular}
\end{table}

\textbf{Önemli bir sınırlılık.} İki normal bilgisayar arasındaki deney gerçek ağ
gecikmesini, seğirmeyi ve bant genişliğini verir. Ancak bu bir \emph{dağıtık
yerleşim} ölçümüdür, bir \emph{sis düğümü} ölçümü değildir: sis tanımındaki
kaynak kısıtı yoktur. Rapor boyunca bu ayrım korunmuştur.

\subsection{Ölçüm kuralları}

Aşağıdaki kurallar bütün koşularda geçerlidir:

\begin{enumerate}
\item \textbf{Yığın boyutu 1.} Gerçek kullanıcı tek pencere gönderir.
\item \textbf{Isınma hariç.} İlk 15 pencere (tek süreçli koşularda 30)
istatistiğe girmez.
\item \textbf{Ortalama, medyan, p95, p99 birlikte.} ``En hızlı çalışma''
raporlanmaz.
\item \textbf{Tekrarlı ve dönüşümlü.} Karşılaştırmalarda üç veya dört tekrar,
sıra A, B, A, B biçimindedir; böylece işlemci ısınmasından kaynaklanan kayma
tüm yapılandırmalara eşit dağılır. Ardışık koşulsaydı en son ölçülen ortam
sistematik olarak yavaş görünürdü.
\item \textbf{Tek iş parçacığı.} Gerekçe hız değil kararlılıktır
(Bölüm~\ref{sec:tekrar}).
\item \textbf{İki ayrı dağılım karıştırılmaz.} Koşu içi dağılım (bir koşudaki
yüzlerce pencerenin p50/p95/p99'u) sistemin en kötü durum davranışını; koşu arası
yayılım (aynı deney tekrarlanınca medyanın oynaması) ölçümün güven aralığını
verir.
\item \textbf{Pencere gecikmesi $\ne$ hesaplama gecikmesi.} 10 saniyelik pencere
ve 5 saniyelik güncelleme aralığı verinin \emph{toplanma} süresidir; bu rapor
\emph{hesaplama} süresini ölçer. İkisi her tabloda ayrı tutulur.
\item \textbf{Her sonuç JSON'a yazılır}; tablolar JSON'dan üretilir, elle sayı
taşınmaz.
\end{enumerate}

\subsection{Koşu envanteri}

Tablo~\ref{tab:kosular} bu raporda kullanılan ölçüm kümelerini listeler.

\begin{table}[t]
\caption{Rapordaki ölçüm kümeleri}
\label{tab:kosular}
\centering
\small
\begin{tabular}{p{1.5cm}p{3.0cm}p{2.6cm}}
\toprule
\textbf{Küme} & \textbf{İçerik} & \textbf{Ölçek}\\
\midrule
Tek süreç & Aşama bazlı profil, \texttt{native} yol & S11+S14, 2123 pencere\\
Kıyas & İplik taraması, ardışık düzen karşılaştırması & S11, 4 tur $\times$ 587 pencere\\
Geri döngü & Yedi metrik, tek ve iki istemci & S11, 10\,dk / S11+S14\\
Sis & Beş XAI durumu, uçtan uca & S11, 600 tik, 44 ölçülen pencere\\
Yerleşim & Uç--sis, dönüşümlü & S11, 3 tur $\times$ 600 tik\\
XAI tarama & 17 yapılandırma, GPU & 40 pencere (duyarlılık: 10)\\
\bottomrule
\end{tabular}
\end{table}

%=======================================================================
\section{Sonuçlar}
%=======================================================================

\subsection{Aşama bazlı hesaplama profili}

Tablo~\ref{tab:asama}, tek süreçli \texttt{native} yolda, nedensel
normalizasyonla ölçülen aşama sürelerini verir. Toplam hesaplama gecikmesi
medyanı 4{,}69\ms, p95 6{,}66\ms'dir. En pahalı iki aşama çıkarım (2{,}68\ms
medyan) ve yeniden örneklemedir (1{,}38\ms medyan).

\begin{table}[t]
\caption{Aşama bazlı hesaplama süresi (S11+S14, \texttt{native} yol,
nedensel normalizasyon, 2123 pencere, 30 ısınma penceresi hariç, tek iplik)}
\label{tab:asama}
\centering
\small
\begin{tabular}{lrrrr}
\toprule
\textbf{Aşama} & \textbf{ort.} & \textbf{medyan} & \textbf{p95} & \textbf{p99}\\
\midrule
Yeniden örnekleme & 1{,}507 & 1{,}378 & 2{,}328 & 3{,}259\\
Ön işleme         & 0{,}251 & 0{,}218 & 0{,}401 & 0{,}924\\
Tensör hazırlığı  & 0{,}159 & 0{,}141 & 0{,}249 & 0{,}513\\
Çıkarım           & 2{,}951 & 2{,}682 & 3{,}997 & 5{,}720\\
Son işleme        & 0{,}108 & 0{,}086 & 0{,}171 & 0{,}380\\
\midrule
\textbf{Toplam}   & \textbf{5{,}034} & \textbf{4{,}689} & \textbf{6{,}658} & \textbf{9{,}170}\\
\bottomrule
\end{tabular}
\\[2pt]
\footnotesize Bütün değerler milisaniyedir. Seğirme (ardışık mutlak fark
ortalaması) 1{,}155\ms; işlem CPU'su ort.\ \%90{,}1; RSS ort.\ 470{,}2\,MB,
tepe 558{,}0\,MB.
\end{table}

Bu profil, sistemin hesap tarafının ne kadar ucuz olduğunun temel kanıtıdır:
$640\times14$'lük bir pencere, tek iş parçacığında 5\ms'nin altında
sınıflandırılmaktadır. Bu sayı, aynı pencerenin \emph{açıklanmasının} maliyetiyle
(Bölüm~\ref{sec:fog-xai}) karşılaştırıldığında iki mertebelik fark ortaya
çıkmaktadır.

\subsection{Ölçümün tekrarlanabilirliği}
\label{sec:tekrar}

Tek bir koşunun sonucu bu makinede güvenilir değildir: aynı yapılandırma arka
arkaya koşulduğunda değerler yukarı kaymaktadır (işlemci frekans ölçeklemesi).
Bu nedenle bütün karşılaştırmalar dönüşümlü tekrarla yapılmıştır.

Tablo~\ref{tab:iplik} iş parçacığı sayısının etkisini gösterir. Medyanlar
arasındaki fark, koşu arası yayılımın altında kalmaktadır; yani bu iş yükünde
iplik sayısı belirleyici değildir. Varsayılanın 1 olmasının gerekçesi hız değil
\textbf{kararlılıktır}: tek iplikte yayılım \%2{,}9 iken iki ve dört iplikte
\%10{,}7--14{,}7'dir. Bu, izole mikro-kıyaslamaların yanıltıcı olduğunu ve
kararın her zaman tam ardışık düzen üzerinde verilmesi gerektiğini gösteren bir
örnektir.

\begin{table}[t]
\caption{İş parçacığı taraması (S11, 4 dönüşümlü tur, koşu içi medyanların
medyanı)}
\label{tab:iplik}
\centering
\small
\begin{tabular}{crrrr}
\toprule
\textbf{İplik} & \textbf{Toplam (ms)} & \textbf{Yayılım} & \textbf{VK} & \textbf{p95 (ms)}\\
\midrule
1 & \textbf{3{,}397} & \textbf{\%2{,}9} & \textbf{\%1{,}2} & 4{,}980\\
2 & 3{,}933 & \%14{,}7 & \%6{,}0 & 5{,}774\\
4 & 3{,}684 & \%10{,}7 & \%4{,}5 & 5{,}508\\
\bottomrule
\end{tabular}
\\[2pt]
\footnotesize VK: değişim katsayısı. Üç yapılandırmada sınıflandırma sonucu
birebir aynıdır (doğruluk 0{,}8586), yani iplik sayısı yalnızca zamanlamayı
etkiler.
\end{table}

\subsection{Yeniden örneklemenin yeri ve bant genişliği}

Tablo~\ref{tab:pipeline}, yeniden örneklemenin ölçüme dahil edilip
edilmemesinin etkisini verir. Yeniden örnekleme tek başına toplamın yaklaşık
\%28{,}6'sıdır, yani çıkarımdan sonraki en pahalı aşamadır. İki yol arasındaki
\%38{,}6'lık fark, koşu arası yayılımın (\%8--9) çok üstünde olduğu için
gerçektir.

\begin{table}[t]
\caption{Ardışık düzen karşılaştırması (S11, 4 dönüşümlü tur, tek iplik)}
\label{tab:pipeline}
\centering
\small
\setlength{\tabcolsep}{4pt}
\begin{tabular}{lrrrrr}
\toprule
\textbf{Yol} & \textbf{Yen.\ ör.} & \textbf{Ön işl.} & \textbf{Çıkarım} & \textbf{Toplam} & \textbf{p95}\\
\midrule
\texttt{prepared} & ---       & 0{,}580 & 2{,}969 & 3{,}971 & 6{,}349\\
\texttt{native}   & 1{,}573   & 0{,}252 & 3{,}039 & 5{,}502 & 7{,}951\\
\bottomrule
\end{tabular}
\\[2pt]
\footnotesize Bütün değerler milisaniyedir. Koşu arası yayılım: \texttt{prepared}
\%9{,}4, \texttt{native} \%8{,}1.
\end{table}

Doğruluk korunmaktadır. Pencere başına yeniden örnekleme, tüm sinyali bir kerede
indirmekle kenar etkisi nedeniyle birebir aynı sonucu vermez, ancak fark
ölçülebilir düzeyin altındadır: \texttt{prepared} 0{,}8586 (587 pencere),
\texttt{native} 0{,}8574 (589 pencere).

\textbf{Bant genişliği karşılığı.} 500\ms'lik bir tikte ham hızda 11\,536 bayt,
64\,Hz'e indirilmiş halde 1\,792 bayt taşınır. Yani indirgemeyi sunucuda yapmak
ağdan \textbf{6{,}4 kat} fazla veri geçirmek demektir. Bu, yerleşim kararının
sayısal dayanağıdır: hesabı sunucuya taşımak ağ maliyetini artırır.

\subsection{Ortak yerleşimde yedi metrik}

Tablo~\ref{tab:loopback}, gerçek sensör hızında, tek istemciyle, S11'in ilk 10
dakikası boyunca geri döngü üzerinden ölçülen yedi metriği verir.

\begin{table}[t]
\caption{Geri döngüde yedi metrik (S11, tek istemci, gerçek zamanlı, 10\,dk,
104 ölçülen pencere)}
\label{tab:loopback}
\centering
\small
\begin{tabular}{lrrrr}
\toprule
\textbf{Metrik} & \textbf{ort.} & \textbf{medyan} & \textbf{p95} & \textbf{p99}\\
\midrule
Gecikme          & 0{,}909 & 0{,}809 & 1{,}466 & 2{,}707\\
Kuyruk gecikmesi & 0{,}062 & 0{,}045 & 0{,}085 & 0{,}622\\
Çalışma süresi   & 8{,}273 & 7{,}179 & 10{,}524 & 11{,}268\\
Yanıt süresi     & 9{,}884 & 8{,}821 & 12{,}698 & 13{,}927\\
\bottomrule
\end{tabular}
\\[2pt]
\footnotesize Değerler milisaniyedir. Seğirme 3{,}529\ms. Bant genişliği
184{,}8\,kbps gönderilen, 2{,}3\,kbps alınan, toplam 187{,}2\,kbps. Sunucu içi
aşama dökümü (ortalama): yeniden örnekleme 2{,}211, ön işleme 0{,}449, tensör
0{,}324, çıkarım 4{,}648, son işleme 0{,}135\ms.
\end{table}

\textbf{Doğrulama.} Ayrı bir geri döngü koşusunda (iki istemci, S11 ve S14'ün
tam kayıtları, 2123 pencere) tek süreçli \texttt{native} yol ile
istemci--sunucu yolu aynı veride birebir aynı sınıflandırma sonucunu
vermektedir: doğruluk 0{,}8222, stres F1 0{,}6944, karışıklık matrisi
TN$=$743, FP$=$184, FN$=$29, TP$=$242. Bu, ağ katmanının hesabı bozmadığının
kanıtıdır. Tablo~\ref{tab:loopback}'teki on dakikalık koşu S11'in yalnızca
başını kapsadığı için sınıflandırma açısından bilgi taşımaz; amacı zamanlama
ölçümüdür.

\textbf{Yakalanan bir hata.} Tek süreçli \texttt{native} yolda, çoğunluk etiketi
geçersiz olan pencereler tamamen atlanıyordu; bu, normalizasyon istatistiğinin o
pencerelerde hiç güncellenmemesi anlamına geliyordu. Gerçek bir sunucu etiketi
görmez, dolayısıyla böyle bir atlama yapamaz. Düzeltildikten sonra doğruluk
0{,}7674'ten 0{,}8222'ye yükselmiş ve iki yol birbirine eşitlenmiştir. Geçersiz
etiketli pencereler artık baştan sona işlenmekte, yalnızca sınıflandırma
raporundan çıkarılmaktadır.

\subsection{Uç--sis yerleşim karşılaştırması}

Tablo~\ref{tab:yerlesim}, aynı deneyin iki yerleşimde, üç dönüşümlü turla
koşulmuş halini verir. Değerler koşu içi medyanların medyanı, parantez içindeki
sayılar koşu arası yayılımdır.

\begin{table}[t]
\caption{Uç (ortak yerleşim) ve sis karşılaştırması (S11, 600 tik, 3 dönüşümlü
tur, açıklayıcı kapalı)}
\label{tab:yerlesim}
\centering
\small
\begin{tabular}{lrr}
\toprule
\textbf{Metrik} & \textbf{Uç (ortak)} & \textbf{Sis (Wi-Fi)}\\
\midrule
Gecikme (medyan)        & 0{,}928 (\%7{,}0)  & 4{,}910 (\%44{,}7)\\
Gecikme (p95)           & 2{,}073            & 66{,}179\\
Kuyruk gecikmesi        & 0{,}048            & 0{,}028\\
Çalışma süresi (medyan) & 6{,}344 (\%15{,}4) & \textbf{3{,}372 (\%6{,}3)}\\
Çalışma süresi (p95)    & 8{,}797            & 4{,}400\\
Yanıt süresi (medyan)   & \textbf{7{,}880}   & 10{,}151\\
Yanıt süresi (p95)      & 10{,}360           & 74{,}355\\
Seğirme                 & 2{,}221            & 22{,}600\\
Bant genişliği (kbps)   & 187{,}17           & 187{,}11\\
Sunucu CPU              & \%0{,}574          & \%0{,}111\\
Sunucu RSS (MB)         & 186{,}8            & 209{,}5\\
Saat belirsizliği       & 0{,}122            & 1{,}359\\
\bottomrule
\end{tabular}
\\[2pt]
\footnotesize Değerler aksi belirtilmedikçe milisaniyedir.
\end{table}

Üç gözlem öne çıkmaktadır. Birincisi, \textbf{sis düğümü hesabı daha hızlı
yapmaktadır} (3{,}372\ms'ye karşı 6{,}344\ms): ortak yerleşimde istemci ve
sunucu aynı işlemciyi paylaşmakta ve birbirlerini yavaşlatmaktadır. İkincisi,
bu kazanç ağ tarafında geri verilmektedir; yanıt süresi medyanı sis lehine
değil, aleyhine çıkmaktadır (10{,}151'e karşı 7{,}880\ms). Üçüncüsü, Wi-Fi'ın
kuyruk davranışı p95'te dramatiktir: yanıt süresi p95'i 74\ms'ye, seğirme
22{,}6\ms'ye çıkmaktadır ve koşu arası yayılım \%112'ye ulaşmaktadır. Bant
genişliği her iki yerleşimde aynıdır, beklendiği gibi.

\subsection{Açıklayıcı hiperparametre taraması}
\label{sec:xai-tarama}

Tarama, sunucunun gördüğü girdi uzayında (nedensel normalizasyon) ve GPU
üzerinde koşulmuştur. \textbf{Kalite ölçütleri donanımdan bağımsızdır ve bu
tablolardan alınır; süre değerleri GPU'ya aittir} ve nihai süre tablosu
Bölüm~\ref{sec:fog-xai}'den gelir. Bu ayrım kritiktir: aynı yapılandırmanın CPU
süresi bir mertebe daha büyüktür.

Taramada tek faktör değiştirme (OFAT) yaklaşımı kullanılmış, her aşamanın
kazananı bir sonraki aşamanın varsayılanı olmuştur. Kalite, 40 pencere üzerinde;
duyarlılık, maliyeti nedeniyle 10 pencere $\times$ 3 deneme üzerinde
hesaplanmıştır.

\begin{table}[t]
\caption{LIME tarama sonuçları (nedensel uzay, GPU, 40 pencere)}
\label{tab:lime-sweep}
\centering
\scriptsize
\begin{tabular}{lrrrrrr}
\toprule
\textbf{Yapılandırma} & \textbf{ort.} & \textbf{p95} & \textbf{Sadakat} & \textbf{Mono.} & \textbf{Duyarl.} & \textbf{RSS}\\
\midrule
\multicolumn{7}{l}{\emph{Aşama 1 --- örnek sayısı (10 dilim, kw 0{,}25)}}\\
250 örnek   & 105{,}3 & 211{,}4 & 0{,}2219 & 0{,}5925 & 0{,}0274 & 1186\\
500 örnek   & 104{,}3 & 162{,}8 & 0{,}2302 & 0{,}6025 & 0{,}0266 & 1203\\
1000 örnek  & 245{,}6 & 376{,}3 & \textbf{0{,}2333} & 0{,}5938 & 0{,}0264 & 1204\\
\midrule
\multicolumn{7}{l}{\emph{Aşama 2 --- bölge sayısı (500 örnek, kw 0{,}25)}}\\
5 dilim     & \textbf{67{,}5} & \textbf{71{,}9} & 0{,}2180 & \textbf{0{,}6175} & \textbf{0{,}0217} & 1205\\
10 dilim    & 151{,}1 & 268{,}1 & 0{,}2302 & 0{,}6025 & 0{,}0266 & 1205\\
20 dilim    & 172{,}9 & 353{,}4 & 0{,}2375 & 0{,}5862 & 0{,}0357 & 1206\\
\midrule
\multicolumn{7}{l}{\emph{Aşama 3 --- çekirdek genişliği (500 örnek, 20 dilim)}}\\
kw 0{,}125  & 104{,}8 & 159{,}4 & 0{,}1887 & \textbf{0{,}6700} & \textbf{0{,}0147} & 1206\\
kw 0{,}25   & 181{,}1 & 358{,}0 & 0{,}2375 & 0{,}5862 & 0{,}0357 & 1206\\
kw 0{,}50   & 152{,}8 & 344{,}5 & \textbf{0{,}2390} & 0{,}5813 & 0{,}0365 & 1206\\
\bottomrule
\end{tabular}
\\[2pt]
\footnotesize Süreler ms, RSS MB. Kendi-kararsızlığı bütün LIME
yapılandırmalarında tam olarak 0{,}0'dır: maske tohumu sabit olduğu için yöntem
deterministiktir ve ölçülen duyarlılık yalnızca girdi gürültüsünden gelir.
\end{table}

\begin{table}[t]
\caption{SHAP tarama sonuçları (nedensel uzay, GPU, 40 pencere,
GradientExplainer)}
\label{tab:shap-sweep}
\centering
\scriptsize
\begin{tabular}{lrrrrrr}
\toprule
\textbf{Yapılandırma} & \textbf{ort.} & \textbf{p95} & \textbf{Sadakat} & \textbf{Mono.} & \textbf{Duyarl.} & \textbf{RSS}\\
\midrule
\multicolumn{7}{l}{\emph{Aşama 1 --- arka plan boyutu ($n{=}200$, rastgele)}}\\
10   & 192{,}3 & 232{,}7 & 0{,}2038 & \textbf{0{,}6263} & \textbf{0{,}1119} & 1407\\
50   & 207{,}1 & 247{,}7 & 0{,}2212 & 0{,}6150 & 0{,}1003 & 1409\\
100  & 204{,}4 & 252{,}8 & 0{,}2243 & 0{,}5938 & 0{,}1011 & 1411\\
\midrule
\multicolumn{7}{l}{\emph{Aşama 2 --- örnek sayısı (arka plan 100, rastgele)}}\\
50   & \textbf{46{,}5} & \textbf{54{,}9} & 0{,}2214 & 0{,}5950 & 0{,}1083 & 1413\\
200  & 204{,}9 & 249{,}4 & 0{,}2243 & 0{,}5938 & 0{,}1011 & 1416\\
400  & 402{,}2 & 487{,}7 & \textbf{0{,}2246} & 0{,}6025 & 0{,}1032 & 1468\\
\midrule
\multicolumn{7}{l}{\emph{Aşama 3 --- arka plan seçimi (100, $n{=}200$)}}\\
rastgele    & 198{,}4 & 234{,}7 & 0{,}2243 & 0{,}5938 & 0{,}1011 & 1413\\
k-ortalama  & 199{,}3 & 240{,}5 & \textbf{0{,}2302} & 0{,}5625 & 0{,}1184 & 1406\\
\bottomrule
\end{tabular}
\\[2pt]
\footnotesize Süreler ms, RSS MB. Kendi-kararsızlığı $\approx 10^{-8}$
mertebesindedir; tohum sabitlemesi sayesinde yöntem pratikte deterministiktir.
\end{table}

Şekil~\ref{fig:pareto} bu iki tabloyu kalite--gecikme düzleminde birleştirir.
Üç bulgu öne çıkmaktadır:

\begin{enumerate}
\item \textbf{Kalite doyuma erken ulaşmaktadır.} LIME'da örnek sayısını 250'den
1000'e çıkarmak sadakati yalnızca \%5{,}1 artırırken gecikmeyi 2{,}3 katına
çıkarmaktadır. SHAP'ta örnek sayısını 50'den 400'e çıkarmak sadakati \%1{,}4
artırırken gecikmeyi 8{,}6 katına çıkarmaktadır. \emph{Pareto cephesi
dikleşmektedir}: ucuz çalışma noktaları neredeyse bedavadır.
\item \textbf{İki yöntemin kararlılığı aynı mertebede değildir.} LIME'ın
duyarlılığı 0{,}015--0{,}036, SHAP'ınki 0{,}100--0{,}118 aralığındadır; yani
girdiye eklenen küçük gürültüye karşı LIME bu kurulumda üç ilâ yedi kat daha
kararlıdır. Bu, bitirme projesindeki çevrimdışı sıralamanın tersidir ve
nedenlerinden biri bölge tabanlı LIME'ın sabit tohumla deterministik hale
getirilmiş olmasıdır.
\item \textbf{Bellek ayak izi yöntem seçimine bağlıdır.} SHAP, arka plan
kümesini bellekte tuttuğu için LIME'dan yaklaşık 200\,MB fazla yer kaplamaktadır
(1406--1468\,MB'ye karşı 1186--1206\,MB, tarama ortamında).
\end{enumerate}

\begin{figure}[t]
\centering
\begin{tikzpicture}
\begin{axis}[
  width=\columnwidth, height=5.6cm,
  xlabel={Açıklama süresi, ortalama (ms)},
  ylabel={Sadakat},
  xmin=0, xmax=430, ymin=0.18, ymax=0.245,
  grid=both, grid style={gray!20},
  legend pos=south east, legend style={font=\scriptsize,draw=gray!50},
  label style={font=\scriptsize}, tick label style={font=\scriptsize},
  scaled y ticks=false, y tick label style={/pgf/number format/fixed,
    /pgf/number format/precision=3}
]
\addplot[only marks,mark=*,mark size=1.7pt,color=blue!70!black] coordinates {
 (105.3,0.2219) (104.3,0.2302) (245.6,0.2333)
 (67.5,0.2180) (151.1,0.2302) (172.9,0.2375)
 (104.8,0.1887) (181.1,0.2375) (152.8,0.2390)
};
\addlegendentry{LIME}
\addplot[only marks,mark=square*,mark size=1.7pt,color=red!70!black] coordinates {
 (192.3,0.2038) (207.1,0.2212) (204.4,0.2243)
 (46.5,0.2214) (204.9,0.2243) (402.2,0.2246)
 (198.4,0.2243) (199.3,0.2302)
};
\addlegendentry{SHAP}
\end{axis}
\end{tikzpicture}
\caption{Açıklama kalitesine karşı açıklama süresi (GPU, nedensel uzay). Sol üst
köşe tercih edilir. Her iki yöntemde de kalite dar bir bantta toplanırken süre
bir mertebe değişmektedir.}
\label{fig:pareto}
\end{figure}

\textbf{Metodolojik bir karşı örnek.} Tarama önce sızıntılı (denek bazlı)
normalizasyon uzayında koşulmuş, sonra nedensel uzayda tekrarlanmıştır. İki
ölçüm yan yana konduğunda sızıntılı ölçümün XAI'yi olduğundan iyi gösterdiği
görülmektedir: LIME'ın sadakati 0{,}2476'dan 0{,}2333'e düşmekte, duyarlılığı
0{,}0174'ten 0{,}0264'e çıkmaktadır. Yani açıklayıcıyı, dağıtımda hiç
görmeyeceği bir girdi dağılımı üzerinde değerlendirmek iyimser bir kalite
tahmini üretir. Bu, XAI değerlendirmesinin sunucunun gerçekten gördüğü girdi
uzayında yapılması gerektiğinin doğrudan kanıtıdır.

Bu taramaya dayanarak dört çalışma noktası dondurulmuş ve sunucuya gömülmüştür
(Tablo~\ref{tab:secilen}). İki maliyet noktası kasıtlıdır: seçici açıklamanın
kazancı, açıklama pahalılaştıkça büyür.

\begin{table}[t]
\caption{Sunucuya gömülen açıklayıcı yapılandırmaları}
\label{tab:secilen}
\centering
\small
\begin{tabular}{llp{3.2cm}}
\toprule
\textbf{Ad} & \textbf{Yöntem} & \textbf{Yapılandırma}\\
\midrule
\texttt{lime}      & LIME & 500 örnek, 20 dilim, kw 0{,}5\\
\texttt{lime\_ucuz}& LIME & 250 örnek, 20 dilim, kw 0{,}5\\
\texttt{shap}      & SHAP & arka plan 100 (k-ortalama), $n{=}200$\\
\texttt{shap\_ucuz}& SHAP & arka plan 100 (k-ortalama), $n{=}50$\\
\bottomrule
\end{tabular}
\end{table}

\subsection{Sis düğümünde uçtan uca açıklama maliyeti}
\label{sec:fog-xai}

Bu, raporun ana sonucudur. Beş durum aynı koşullarda (S11, tek istemci, 600 tik,
Wi-Fi, 44 ölçülen pencere, saat farkı 60 saniyede bir yeniden ölçülerek) uçtan
uca koşulmuştur.

\begin{table*}[t]
\caption{Sis düğümünde beş XAI durumunun uçtan uca ölçümü (S11, tek istemci,
600 tik, Wi-Fi, 44 ölçülen pencere, her koşu $\approx$5 dakika)}
\label{tab:fog-xai}
\centering
\small
\begin{tabular}{lrrrrrrrr}
\toprule
\textbf{Açıklayıcı} & \textbf{Yanıt} & \textbf{Yanıt} & \textbf{Çalışma} & \textbf{Açıklama} & \textbf{Açıklama} & \textbf{Gecikme} & \textbf{Sunucu} & \textbf{Sunucu}\\
 & \textbf{medyan} & \textbf{p95} & \textbf{medyan} & \textbf{medyan} & \textbf{p95} & \textbf{medyan} & \textbf{CPU} & \textbf{RSS tepe}\\
\midrule
Yok                          & \textbf{10{,}6}   & 42{,}2    & 3{,}5     & ---       & ---       & 4{,}5 & \%0{,}0  & 210{,}4\\
SHAP ($n{=}50$)              & 442{,}6           & 574{,}0   & 394{,}1   & 390{,}4   & 460{,}5   & 4{,}3 & \%5{,}0  & 422{,}3\\
LIME (250 örnek)             & 645{,}2           & 802{,}5   & 588{,}1   & 584{,}0   & 650{,}6   & 5{,}3 & \%10{,}8 & 442{,}4\\
LIME (500 örnek)             & 1227{,}5          & 1407{,}1  & 1158{,}1  & 1154{,}2  & 1295{,}9  & 5{,}0 & \%22{,}8 & 449{,}2\\
SHAP ($n{=}200$)             & \textbf{1456{,}2} & 1548{,}5  & 1431{,}1  & 1426{,}2  & 1465{,}0  & 4{,}5 & \%25{,}9 & 464{,}5\\
\bottomrule
\end{tabular}
\\[2pt]
\footnotesize Süreler milisaniye, RSS megabayttır. CPU değeri sunucu sürecinin
koşu boyunca ortalamasıdır. Bant genişliği beş durumda da 186{,}7--187{,}2\,kbps
aralığındadır. Kuyruk gecikmesi medyanı 0{,}029--0{,}030\ms'dir.
\end{table*}

Tablo~\ref{tab:fog-xai} ve Şekil~\ref{fig:fogbar} üç sonuç vermektedir:

\begin{enumerate}
\item \textbf{Açıklama, sistemin baskın maliyetidir.} Açıklayıcı kapalıyken
yanıt süresi medyanı 10{,}6\ms'dir; en pahalı yapılandırmada 1456{,}2\ms,
yani \textbf{137 kat}. Sunucu içi aşama dökümünde açıklama, çalışma süresinin
\%99{,}6'sını oluşturmaktadır (1426{,}2 / 1431{,}1\ms). Çıkarımın kendisi hâlâ
2--3{,}4\ms'dir; yani modelin maliyeti açıklamanın yanında ölçüm gürültüsü
mertebesindedir.

\item \textbf{Ucuz çalışma noktaları gerçek zamanlılık sınırının altındadır.}
ShaTS'in \cite{shats2025} koyduğu 1\,s/pencere sınırına göre SHAP ($n{=}50$,
442{,}6\ms) ve LIME (250 örnek, 645{,}2\ms) uygundur; pahalı yapılandırmalar
medyanda ve p95'te bu sınırı aşmaktadır. Beş saniyelik güncelleme aralığı
bütçesine göre bakıldığında ucuz noktalar bütçenin \%9--13'ünü, pahalı noktalar
\%25--29'unu tüketmektedir --- tek kullanıcıda. Dört kullanıcıda pahalı
yapılandırma bütçeyi aşacaktır; seçici açıklamanın gerekçesi tam buradadır.

\item \textbf{Bellek, kısıtlı donanımda bağlayıcı olabilir.} RSS tepe değeri
210\,MB'den 464\,MB'ye çıkmaktadır. 2\,GB bellekli bir Raspberry Pi'de bu
sıkışacaktır; en az 4\,GB gerekmektedir. Artışın kaynağı ikilidir: SHAP'ta
bellekte tutulan arka plan kümesi, LIME'da yüzlerce perturbasyonun tek partide
oluşturulması.
\end{enumerate}

\begin{figure}[t]
\centering
\begin{tikzpicture}
\begin{axis}[
  width=\columnwidth, height=5.2cm,
  ybar, bar width=12pt,
  ylabel={Yanıt süresi medyanı (ms)},
  symbolic x coords={Yok,SHAP-50,LIME-250,LIME-500,SHAP-200},
  xtick=data, x tick label style={font=\scriptsize,rotate=20,anchor=north east},
  ymin=0, ymax=1700,
  nodes near coords, nodes near coords style={font=\scriptsize},
  label style={font=\scriptsize}, tick label style={font=\scriptsize},
  grid=major, grid style={gray!20},
  enlarge x limits=0.13
]
\addplot[fill=blue!35,draw=blue!60!black] coordinates {
 (Yok,10.6) (SHAP-50,442.6) (LIME-250,645.2) (LIME-500,1227.5) (SHAP-200,1456.2)
};
\end{axis}
\end{tikzpicture}
\caption{Sis düğümünde yanıt süresi medyanı. Açıklayıcı kapalıyken 10{,}6\ms
olan yanıt, en pahalı yapılandırmada 137 katına çıkmaktadır.}
\label{fig:fogbar}
\end{figure}

\textbf{Tutarlılık denetimi.} Gecikme beş koşuda da 4{,}3--5{,}3\ms bandında
kalmaktadır; yani açıklayıcı seçimi ağ metriklerini etkilememektedir. Bu
beklenen davranıştır ve ölçümün kendi içinde tutarlı olduğunun göstergesidir:
açıklayıcı sunucu tarafında çalıştığı için taşıma katmanına dokunmaz.

\textbf{Kayıt boşluğu ve düzeltmesi.} Bu koşular alındığında hangi açıklayıcının
kullanıldığı yalnızca dosya adında saklanıyordu; sunucu bu bilgiyi istemciye
göndermiyordu. Bu, sonuçların yeniden üretilebilirliği açısından bir eksikti ve
sonradan giderilmiştir: sunucu artık kapanış istatistiklerinde yöntem adını,
eşiği, açıklayıcı yapılandırmasını ve ısınma süresini de göndermektedir.

\subsection{Sınıflandırma başarımının ölçüm yolundan bağımsızlığı}

Tablo~\ref{tab:siniflama}, aynı modelin farklı ölçüm yollarındaki sınıflandırma
sonucunu toplar. İki gözlem önemlidir.

\begin{table}[t]
\caption{Sınıflandırma başarımı, ölçüm yoluna göre}
\label{tab:siniflama}
\centering
\small
\begin{tabular}{lrrr}
\toprule
\textbf{Koşul} & \textbf{Doğruluk} & \textbf{Stres F1} & \textbf{Pencere}\\
\midrule
Çevrimdışı, denek bazlı & 0{,}9361 & --- & ---\\
Tek süreç, \texttt{running}    & 0{,}8222 & 0{,}6944 & 1198\\
İstemci--sunucu, \texttt{running} & 0{,}8222 & 0{,}6944 & 1198\\
Tek süreç, \texttt{fixed}$^{\dagger}$ & 0{,}6996 & 0{,}4676 & 1175\\
\bottomrule
\end{tabular}
\\[2pt]
\footnotesize $^{\dagger}$ Bu satır, yeniden örneklemenin ölçüm dışında
bırakıldığı ve geçersiz etiketli pencerelerin atlandığı önceki
(\texttt{prepared}) ardışık düzen döneminden gelmektedir; \texttt{native} yolda
\texttt{fixed} normalizasyonla bir koşu henüz yapılmamıştır. Diğer satırlarla
birebir karşılaştırılmamalı, yalnızca kişiselleştirme etkisinin mertebesi
olarak okunmalıdır.
\end{table}

Birincisi, \textbf{ağ katmanı hesabı bozmamaktadır}: tek süreçli ve
istemci--sunucu yolları birebir aynı karışıklık matrisini vermektedir.

İkincisi, \textbf{düşüş modelin değil normalizasyon kısıtının sonucudur}.
Çevrimdışı 0{,}9361 değeri, deneğin tüm kaydının istatistiğiyle, yani geleceğe
sızarak elde edilmiştir. Nedensel normalizasyonla 0{,}8222'ye, deneğe hiç
uyarlanmayan sabit istatistikle 0{,}6996'ya düşmektedir. Aradaki 12 puanlık fark,
kişiselleştirmenin gerçek zamanlı bir sistemde ne kadar değer taşıdığının
doğrudan ölçüsüdür.

\subsection{Mikrodenetleyici fizibilitesi}

Açıklama maliyeti sunucuda ölçüldükten sonra doğal bir soru kalmaktadır: çıkarım
giyilebilir cihazın içinde yapılabilir mi? Yapılabilirse uç ile sis arasında
seçim yapma sorunu ortadan kalkar. Bu, donanım olmadan katman katman parametre,
aktivasyon belleği, tepe RAM ihtiyacı ve çarpma-toplama (MAC) sayısı
hesaplanarak analiz edilmiştir.

\begin{table}[t]
\caption{Mikrodenetleyici gereksinimleri (10\,s pencere, 49\,890 parametre,
pencere başına 14{,}50 milyon MAC)}
\label{tab:mcu}
\centering
\small
\begin{tabular}{lrrr}
\toprule
\textbf{Veri tipi} & \textbf{Ağırlık (Flash)} & \textbf{Tepe RAM} & \textbf{Giriş tamponu}\\
\midrule
float32 & 194{,}9\,KB & 320{,}0\,KB & 35{,}0\,KB\\
float16 & 97{,}4\,KB  & 160{,}0\,KB & 17{,}5\,KB\\
int8    & 48{,}7\,KB  & 80{,}0\,KB  & 8{,}8\,KB\\
\bottomrule
\end{tabular}
\\[2pt]
\footnotesize Tepe aktivasyon ikinci evrişim bloğunun BatchNorm katmanındadır
(81\,920 eleman): 128 kanal $\times$ 320 uzunluk.
\end{table}

Sonuç nettir: \textbf{darboğaz hız değil bellektir}. Pencere başına 14{,}50
milyon MAC gerekirken bütçe 5000\ms'dir (güncelleme aralığı). En yavaş senaryo
olan Cortex-M4 @168\,MHz float32'de tahmini süre 287{,}7\ms, yani bütçenin
\%5{,}8'idir; Cortex-M7 @480\,MHz int8'de 12{,}1\ms, yani \%0{,}2'dir. Buna
karşılık float32 ağırlıklar yalnızca STM32H743 ve U575 sınıfı cihazlara,
int8 ağırlıklar ise F407 ve üzerine sığmaktadır; F103 sınıfı hiçbir biçimde
yetmemektedir.

Bu analiz kâğıt üstündedir. Cihaza indirmek için ONNX aktarımı, STM32Cube.AI ya
da TensorFlow Lite Micro ile kod üretimi, int8 nicemleme ve doğruluk kaybı
ölçümü, ön işlemenin (700\,Hz'den 64\,Hz'e indirgeme ve nedensel z-skor) porte
edilmesi ve cihaz üstü zaman ölçümü gerekmektedir.

%=======================================================================
\section{Tartışma}
%=======================================================================

\textbf{Açıklama maliyeti, model maliyetinin iki mertebe üzerindedir.} Bu,
raporun en net bulgusudur. 49\,890 parametreli bir modelin çıkarımı 2--3\ms
sürerken, aynı pencerenin açıklanması 390--1426\ms sürmektedir. Literatürdeki
``gerçek zamanlı kullanıma uygundur'' iddialarının büyük çoğunluğu yalnızca
çıkarım süresine bakmaktadır; açıklama katmanı eklendiğinde bu iddialar test
edilmemiş varsayımlara dönüşmektedir.

\textbf{Kalite doyuma erken ulaşmakta, maliyet ulaşmamaktadır.} Tarama
tablolarında sadakat değerleri 0{,}19--0{,}24 gibi dar bir banttayken süre
46{,}5\ms ile 402{,}2\ms arasında değişmektedir. Bu asimetri, ucuz çalışma
noktalarının pratikte neredeyse bedava olduğu anlamına gelir ve hiperparametre
taramasının \emph{süreyle birlikte} yapılmasının niçin gerektiğini gösterir.
Yalnızca kaliteye bakan otomatik bir seçim, bu taramada her aşamada en pahalı
yapılandırmayı seçerdi; bu nedenle nihai seçim, kalite--gecikme tablosuna
bakılarak elle yapılmıştır.

\textbf{Çevrimdışı sıralama gerçek zamanda korunmayabilir.} Bitirme projesinde
SHAP en kararlı yöntemdi (duyarlılık 0{,}0038'e karşı LIME 0{,}0116). Bu
kurulumda, nedensel uzayda ve bölge tabanlı LIME ile sıralama tersine dönmüştür
(LIME 0{,}015--0{,}036, SHAP 0{,}100--0{,}118). İki ölçüm farklı modeller,
farklı girdi uzayları ve farklı LIME uygulamaları üzerinde yapıldığı için
doğrudan karşılaştırılamaz; ancak bulgu, ``çevrimdışı en iyi olan yöntemin
gerçek zamanda da en iyi olacağı'' varsayımının sınanması gerektiğini
göstermektedir.

\textbf{Yerleşim kararı tek bir metrikle verilemez.} Sis düğümü hesabı daha hızlı
yapmakta (3{,}4'e karşı 6{,}3\ms) ama Wi-Fi gecikmesi ve seğirmesi bu kazancı
geri almaktadır. Aynı zamanda ham sinyali sunucuya göndermek, indirgenmiş
sinyale göre 6{,}4 kat bant genişliği maliyeti getirmektedir. Karar, hangi
metriğin uygulamada bağlayıcı olduğuna bağlıdır.

\textbf{Ölçüm altyapısının kendisi bir katkıdır.} Bu çalışmada üç ayrı ölçüm
hatası yakalanmış ve düzeltilmiştir: (i) 15{,}6\ms çözünürlüklü saatin gecikmeyi
sıfır okuması, (ii) kuyruğun olmaması nedeniyle kuyruk gecikmesinin gecikme
içinde gizlenmesi, (iii) etiketi geçersiz pencerelerin atlanması nedeniyle
normalizasyon istatistiğinin güncellenmemesi. Üçü de sessiz hatalardı; sonuçlar
makul görünüyor ama yanlıştı. Gerçek zamanlı ölçüm yapan çalışmaların bu tür
denetimleri açıkça raporlaması gerektiğini düşünüyoruz.

%=======================================================================
\section{Sınırlılıklar}
%=======================================================================

\begin{enumerate}
\item \textbf{Emülasyon, canlı sistem değil.} Sinyal gerçek sensörden değil,
WESAD kaydından gerçek hızda yeniden oynatılmaktadır. Referans çalışma
\cite{erzurumluoglu2025} da DREAMER'ı aynı biçimde akıtmıştır; yöntemin
savunulabilirliği buna dayanır.

\item \textbf{Sis düğümü kaynak kısıtlı değildir.} İkinci makine 12 çekirdekli
bir masaüstü sınıfı işlemcidir. Ölçülen şey \emph{dağıtık yerleşim}dir,
\emph{kaynak kısıtlı sis düğümü} değil. Kısıtlı mod (tek çekirdeğe kilitleme +
tek iplik) altyapıda hazırdır ancak bu rapordaki sayılar kısıtsız modda
alınmıştır.

\item \textbf{İki tutulmuş denek.} S11 ve S14 dışındaki denekler modelin
eğitiminde kullanılmıştır. Eşik ve kalibrasyon çalışmaları için çapraz denek
protokolü (S11'de kalibre et, S14'te ölç ve tersi) tasarlanmış fakat henüz
koşulmamıştır.

\item \textbf{Sis koşularındaki sınıflandırma sayıları anlamlı değildir.}
Tablo~\ref{tab:fog-xai}'deki koşular S11'in ilk beş dakikasını kapsamaktadır ve
bu aralıkta stres etiketi hiç bulunmamaktadır (42 etiketli pencerenin tamamı
stres dışı). Bu koşuların amacı zamanlama ölçümüdür; sınıflandırma başarımı
Tablo~\ref{tab:siniflama}'den okunmalıdır.

\item \textbf{Bulut ayağı yoktur.} Uç--sis karşılaştırması yapılmış, bulut
yerleşimi bütçe nedeniyle atlanmıştır. Altyapı bunu tek bir adres parametresiyle
desteklemektedir.

\item \textbf{Wi-Fi'da gecikme belirsizliği yüksektir.} Tek yönlü gecikmenin hata
payı $\pm1$--$10$\ms'dir ve ölçülen değerle aynı mertebeye çıkabilir. Çalışma
süresi, kuyruk gecikmesi, yanıt süresi ve seğirme bundan etkilenmez; gecikme
sütunu dikkatle okunmalıdır. Ethernet üzerinde tekrar yapılmamıştır.

\item \textbf{Referans çalışmanın birincil kaynağı doğrulanamamıştır.}
\cite{erzurumluoglu2025}'in erişilebilir bir kaydı bulunamamıştır; bu çalışmayla
yapılan metrik karşılaştırması, künye ve sayılar birincil kaynaktan
doğrulanmadan nihai metne alınmamalıdır.

\item \textbf{Tek model.} Açıklama maliyeti yalnızca TemporalCNN üzerinde
ölçülmüştür. TranSensFusers gibi daha büyük modellerde hem çıkarım hem de
gradyan tabanlı açıklama maliyeti artacaktır; oran değişebilir.
\end{enumerate}

%=======================================================================
\section{Devam Eden İşler}
\label{sec:devam}
%=======================================================================

Seçici açıklama mekanizması sunucuya entegre edilmiş, çalıştığı doğrulanmış ve
iki uç durumu ölçülmüştür. Aşağıdaki adımlar planlanmış fakat henüz
tamamlanmamıştır:

\begin{enumerate}
\item \textbf{Güven kalibrasyonu.} Güvenilirlik diyagramı ve beklenen
kalibrasyon hatası (ECE); gerekirse tek parametreli sıcaklık ölçekleme.
Kalibrasyon, kalibrasyon deneğinde öğrenilip test deneğinde uygulanacaktır.

\item \textbf{Eşik taraması.} $\tau \in \{0{,}60,\ 0{,}70,\ 0{,}80,\ 0{,}90\}$
için: açıklanan pencere yüzdesi, yakalanan hata yüzdesi (yanlış sınıflanan
pencerelerin kaçında açıklama tetiklendi), pencere başına ortalama hesap, p95
gecikme, CPU ve RSS. Bu, ``belirli bir açıklama kapsamında hataların ne kadarı
açıklandı'' çalışma eğrisini verecektir.

\item \textbf{Çok istemcili yük.} 1, 2 ve 4 istemciyle kuyruk gecikmesi p95'inin
büyümesi. Seçici açıklamanın kazancının asıl görüneceği yer burasıdır: dört
kullanıcı $\times$ 400\ms, beş saniyelik bütçenin üçte birini tüketmektedir.

\item \textbf{Düşük güvenli pencerelerde açıklama kararsızlığı.} En güçlü
itiraza cevap: post-hoc açıklamalar karar sınırına yakın oynaktır ve kapı tam
orada tetiklenmektedir. Duyarlılık ölçütü yalnızca düşük güvenli pencerelerde
ayrıca hesaplanacaktır.

\item \textbf{Model karmaşıklığı ablasyonu.} Mevcut OFAT taramasından derinlik,
filtre ve çekirdek varyantlarının F1--gecikme Pareto eğrisi.
\end{enumerate}

%=======================================================================
\section{Sonuç}
%=======================================================================

Bu rapor, çevrimdışı olarak seçilmiş bir fizyolojik stres sınıflandırıcısını
gerçek zamanlı bir istemci--sunucu sistemine yerleştirmiş ve açıklanabilirliğin
hesaplama maliyetini, taşıma katmanı metrikleriyle aynı bütçede ölçmüştür.

Ana bulgu, açıklama maliyetinin model maliyetini iki mertebe aşmasıdır: ayrı bir
makinede, Wi-Fi üzerinden ölçülen yanıt süresi medyanı açıklayıcı kapalıyken
10{,}6\ms, en pahalı yapılandırmada 1456{,}2\ms'dir. Sunucu bellek ayak izi
210\,MB'den 464\,MB'ye çıkmaktadır. Hiperparametre taraması, açıklama kalitesinin
dar bir bantta doyuma ulaştığını fakat sürenin bir mertebe değiştiğini
göstermiştir; bu asimetri, ucuz çalışma noktalarının (SHAP 442{,}6\ms, LIME
645{,}2\ms) literatürdeki bir saniyelik gerçek zamanlılık sınırının altında
kalmasını sağlamaktadır.

Bu sayılar, güven eşikli seçici açıklamanın niçin gerekli olduğunu niceliksel
olarak ortaya koymaktadır: tek kullanıcıda bile pahalı bir açıklayıcı güncelleme
aralığı bütçesinin dörtte birini tüketmekte, çok kullanıcılı ya da kaynak
kısıtlı senaryolarda ise bütçeyi aşmaktadır. Mekanizma sisteme entegre edilmiş ve
uç durumları ölçülmüştür; eşik taraması ve kalibrasyon çalışmaları devam
etmektedir.

%=======================================================================
\begin{thebibliography}{99}
\small

\bibitem{schmidt2018}
P. Schmidt, A. Reiss, R. Duerichen, C. Marberger ve K. Van Laerhoven,
``Introducing WESAD, a multimodal dataset for wearable stress and affect
detection,'' \emph{Proc. 20th ACM Int. Conf. Multimodal Interaction}, 2018,
ss. 400--408.

\bibitem{kasnesis2025}
P. Kasnesis, C. Chatzigeorgiou, M. Feidakis, A. Gutierrez ve C. Z. Patrikakis,
``TranSenseFusers: A temporal CNN-transformer neural network family for
explainable PPG-based stress detection,'' \emph{Biomedical Signal Processing
and Control}, c. 102, s. 107248, 2025.

\bibitem{ribeiro2016}
M. T. Ribeiro, S. Singh ve C. Guestrin, ```Why should I trust you?' Explaining
the predictions of any classifier,'' \emph{Proc. 22nd ACM SIGKDD}, 2016,
ss. 1135--1144.

\bibitem{lundberg2017}
S. M. Lundberg ve S.-I. Lee, ``A unified approach to interpreting model
predictions,'' \emph{Advances in Neural Information Processing Systems}, c. 30,
2017.

\bibitem{sundararajan2017}
M. Sundararajan, A. Taly ve Q. Yan, ``Axiomatic attribution for deep networks,''
\emph{Proc. Int. Conf. Machine Learning (ICML)}, 2017, ss. 3319--3328.

\bibitem{lin2017}
T.-Y. Lin, P. Goyal, R. Girshick, K. He ve P. Dollár, ``Focal loss for dense
object detection,'' \emph{Proc. IEEE ICCV}, 2017, ss. 2980--2988.

\bibitem{nauta2023}
M. Nauta \emph{ve ark.}, ``From anecdotal evidence to quantitative evaluation
methods: A systematic review on evaluating explainable AI,'' \emph{ACM Computing
Surveys}, c. 55, sy. 13s, 2023.

\bibitem{shats2025}
M. Franco de la Peña \emph{ve ark.}, ``ShaTS: A Shapley-based method for
time-series explainability in industrial IoT,'' arXiv:2506.01450, 2025.

\bibitem{erzurumluoglu2025}
E. Erzurumluoğlu \emph{ve ark.}, ``Emotion recognition on a fog computing
architecture,'' \emph{Proc. Int. Conf. Advanced Communication Technology
(ICACT)}, 2025.

\bibitem{basaran2025}
M. Basaran ve F. Dressler, ``XAI-on-RAN: Explainability overhead in live O-RAN
control loops,'' arXiv:2511.17514, 2025.

\bibitem{taguta2026}
T. Taguta \emph{ve ark.}, ``Conditional SHAP for on-device cold-chain
monitoring: A 24-day field deployment,'' arXiv:2609.14036, 2026.

\bibitem{das2021}
D. Das \emph{ve ark.}, ``Explainable activity recognition from wearable
sensors,'' \emph{ACM Trans. Interactive Intelligent Systems}, 2021;
arXiv:2105.09787.

\bibitem{warnecke2020}
A. Warnecke, D. Arp, C. Wressnegger ve K. Rieck, ``Evaluating explanation
methods for deep learning in security,'' \emph{Proc. IEEE European Symp. on
Security and Privacy (EuroS\&P)}, 2020.

\bibitem{jethani2022}
N. Jethani, M. Sudarshan, I. Covert, S.-I. Lee ve R. Ranganath, ``FastSHAP:
Real-time Shapley value estimation,'' \emph{Proc. Int. Conf. Learning
Representations (ICLR)}, 2022.

\bibitem{almadhor2025}
A. Almadhor \emph{ve ark.}, ``Explainable deep learning for multimodal
physiological stress and emotion recognition,'' \emph{Frontiers in
Bioengineering and Biotechnology}, 2025.

\bibitem{elrashidy2025}
N. El Rashidy \emph{ve ark.}, ``An explainable fog-based emotion recognition
framework,'' \emph{Frontiers in Computer Science}, c. 7, 2025.

\bibitem{paes2024}
L. M. Paes, D. Wei ve F. P. Calmon, ``Selective explanations,''
arXiv:2405.19562, 2024.

\bibitem{teerapittayanon2017}
S. Teerapittayanon, B. McDanel ve H. T. Kung, ``BranchyNet: Fast inference via
early exiting from deep neural networks,'' arXiv:1709.01686, 2017.

\bibitem{berestizshevsky2018}
K. Berestizshevsky ve G. Even, ``Dynamically sacrificing accuracy for reduced
computation: Cascaded inference based on softmax confidence,''
arXiv:1805.10982, 2018.

\bibitem{artelt2023}
A. Artelt, R. Visser ve B. Hammer, ``Model agnostic local explanations of reject,''
\emph{Neurocomputing}, c. 558, 2023; arXiv:2205.07623.

\bibitem{bitirme2026}
S. Cindioğlu ve M. Eldemery, ``Explainable artificial intelligence for stress
detection using physiological data,'' Bitirme Projesi, Yıldız Teknik
Üniversitesi, Bilgisayar Mühendisliği Bölümü, Ocak 2026.

\end{thebibliography}

\end{document}
